% =====================================================================
%  DLoc Compression — Figures Explained (Results & Discussion -> end)
%  Self-contained: put the figure PDFs in the SAME folder and compile
%  with pdfLaTeX (Overleaf-ready).
% =====================================================================
\documentclass[11pt]{article}

\usepackage[a4paper,margin=2.2cm]{geometry}
\usepackage{amsmath,amssymb}
\usepackage{graphicx}
\usepackage{booktabs}
\usepackage{xcolor}
\usepackage{enumitem}
\usepackage{tcolorbox}
\tcbuselibrary{skins,breakable}
\usepackage[hidelinks]{hyperref}

\newtcolorbox{showbox}[1][What it shows (simple words)]{%
  colback=blue!4, colframe=blue!45!black, fonttitle=\bfseries,
  title=#1, breakable, left=2mm, right=2mm, top=1mm, bottom=1mm}
\newtcolorbox{measurebox}[1][What it measures exactly]{%
  colback=purple!4, colframe=purple!50!black, fonttitle=\bfseries,
  title=#1, breakable, left=2mm, right=2mm, top=1mm, bottom=1mm}
\newtcolorbox{madebox}[1][How we made it]{%
  colback=gray!6, colframe=gray!55!black, fonttitle=\bfseries,
  title=#1, breakable, left=2mm, right=2mm, top=1mm, bottom=1mm}
\newtcolorbox{readbox}[1][What to read from it]{%
  colback=green!5, colframe=green!45!black, fonttitle=\bfseries,
  title=#1, breakable, left=2mm, right=2mm, top=1mm, bottom=1mm}

\newcommand{\figpdf}[2][0.8]{\begin{center}\includegraphics[width=#1\linewidth]{#2}\end{center}}

\title{\textbf{DLoc Compression --- Every Figure Explained}\\[2mm]
\large Results \& Discussion to the End of the Thesis:\\
what each figure shows, what it measures, and how we made it}
\date{}
\begin{document}
\maketitle
\tableofcontents
\vspace{3mm}\hrule\vspace{3mm}

% =====================================================================
\section{First: how to read these figures}
% =====================================================================

Every result figure is built from one basic quantity: the \textbf{localization error of each test
sample}. Understand that and the metric words (median, P90, CDF) become easy.

\subsection*{Localization error (the atom)}
For one test sample, the model outputs a heatmap. We take the \textbf{brightest pixel} (the
\emph{argmax}) as the \textbf{predicted location}, and compare it to the \textbf{true location} (the
center of the ground-truth Gaussian bump). The straight-line (Euclidean) distance between them, in
\textbf{meters}, is that sample's error. With $N = 4{,}260$ in-distribution test samples we get 4{,}260
error values; the figures summarize those numbers in different ways.

\subsection*{The summary numbers}
\begin{itemize}[leftmargin=*,noitemsep]
  \item \textbf{Median} --- sort all per-sample errors and take the middle one. ``Half the predictions
        are better than this.'' Robust to a few wild outliers; the headline accuracy number.
  \item \textbf{P90 / P99} --- the 90th / 99th percentile. ``90\% (or 99\%) of predictions are within
        this distance.'' These describe the \emph{tail} --- the bad cases.
  \item \textbf{Mean} --- the average; gets pulled upward by outliers, so it is usually larger than the
        median.
\end{itemize}

\subsection*{The CDF (cumulative distribution function)}
A CDF is a curve. The \textbf{x-axis is error in meters}; the \textbf{y-axis is the fraction of samples
with error $\le x$}. So if the curve passes through $(1\,\text{m},\,0.7)$, then \textbf{70\% of
predictions are within 1\,m}. The \textbf{median} is where the curve crosses $y = 0.5$; \textbf{P90}
where it crosses $y = 0.9$.

\begin{readbox}[How to read any CDF in this thesis]
\textbf{Higher and more to the left = better} (more predictions achieve low error). A curve that climbs
slowly / stays low means a model that rarely gets close. When two curves cross, the better model is the
one that is higher at the error you care about.
\end{readbox}

\subsection*{The qualitative heatmaps}
Some figures show the \emph{actual} predicted heatmaps for a few example samples. In every such panel:
the colored blob is the model's predicted probability map (brightened so the sparse bump is visible), a
\textbf{red cross ($\times$)} marks the \textbf{predicted} peak, and a \textbf{white circle ($\circ$)}
marks the \textbf{true} location. When $\times$ and $\circ$ sit on top of each other, the prediction is
good; when they are far apart, that is an error case.

\begin{measurebox}[An honesty note on the median]
The error is computed on a $0.1$\,m pixel grid, so the median is \emph{quantized} to a lattice of about
$0.1\sqrt{n}$ meters. Two medians within $\sim$0.1\,m are effectively equal --- which is why we also
report mean, CDF, and P90/P99 rather than the median alone.
\end{measurebox}

% =====================================================================
\section{How the figures were generated (the pipeline)}
% =====================================================================

All figures are produced by scripts in \texttt{code/figures/} from the \textbf{released weights}
(HuggingFace \texttt{JustAGeek/dloc-code}) and the \textbf{released data} (HuggingFace
\texttt{JustAGeek/dloc-wild-fig10b}). They are \emph{reproducible}, not drawn by hand.

\begin{enumerate}[leftmargin=*]
  \item \texttt{dump\_results.py} --- the heavy step. For each model/test set it loads the trained
        weights, runs every test sample through the model, and (a) computes the \textbf{per-sample
        error} array (\texttt{errors\_<job>.npy}) and (b) saves six \textbf{representative predicted
        heatmaps} (\texttt{heatmaps\_<job>.npz}), chosen at the error percentiles 2, 25, 50, 75, 90, 98
        (so you see a best case, typical cases, and a worst case).
  \item \texttt{plot\_cdfs.py} --- turns the \texttt{errors\_*.npy} arrays into the CDF curves
        (in-distribution and cross-session).
  \item \texttt{plot\_heatmaps.py} --- renders the six-panel qualitative figures from
        \texttt{heatmaps\_*.npz} (red $\times$ = predicted peak, white $\circ$ = truth; a
        \texttt{PowerNorm} with $\gamma = 0.4$ brightens the otherwise-faint Gaussian).
  \item \texttt{plot\_crosssession.py} --- draws the leave-one-session-out median bars (our three-seed
        mean$\pm$std vs.\ the published DLoc) from \texttt{code/results\_crosssession/*.json}.
  \item \texttt{plot\_compression.py} --- draws the accuracy-vs-size and parameter-count plots from the
        verified efficiency numbers.
\end{enumerate}

\begin{madebox}[Reproducibility detail worth knowing]
The error is a per-sample \emph{argmax} loop. On a many-core machine, NumPy's threaded backend spins its
whole thread pool for each tiny operation and the job appears to hang. The fix is to cap the math
libraries to one thread (\texttt{OMP\_NUM\_THREADS=1}, etc.). With that cap, the regenerated medians
match the thesis exactly (U-Net 0.640, MobileNetV2 0.707, TinyCNN 3.91), confirming the figures are
faithful.
\end{madebox}

% =====================================================================
\section{In-Distribution Figures (same-environment accuracy)}
% =====================================================================

These come from training and testing in the \textbf{same July environment} (17{,}574 train / 4{,}260
test). They answer: \emph{``in the easy case, how accurate is each model?''}

% ---------------------------------------------------------------
\subsection{In-distribution error CDF (Thesis Fig.~5.1)}
\figpdf[0.8]{cdf_fig10b.pdf}

\begin{showbox}
The error curves of the three convolutional students on the same July test set. It compares, across the
\emph{whole} range of errors (not just one number), how often each model gets close.
\end{showbox}

\begin{measurebox}
x-axis: localization error in meters. y-axis: fraction of the 4{,}260 test samples with error $\le x$.
The height of each curve at a given x is ``what fraction of predictions are at least that good.''
\end{measurebox}

\begin{madebox}
\texttt{dump\_results.py} produced \texttt{errors\_unet\_fig10b.npy},
\texttt{errors\_mobilenet\_fig10b.npy}, \texttt{errors\_tinycnn\_fig10b.npy} (one error per test
sample); \texttt{plot\_cdfs.py} sorted each array and plotted the cumulative fraction.
\end{madebox}

\begin{readbox}
The \textbf{MobileNetV2-UNet} curve sits \emph{above} the plain \textbf{MobileNetV2} across the whole
range (0.64\,m vs 0.71\,m median --- consistently better, not just on average). \textbf{TinyCNN} trails
badly --- it never even reaches the 0.8 level within 3\,m, the visual signature of a model below the
capacity floor.
\end{readbox}

% ---------------------------------------------------------------
\subsection{Qualitative in-distribution predictions --- U-Net (Thesis Fig.~5.2)}
\figpdf[1.0]{heatmaps_unet_fig10b.pdf}

\begin{showbox}
Six \emph{actual} predicted heatmaps from the U-Net, picked to span the error range (from a near-perfect
case to a worst case). This shows \emph{what} a good vs.\ bad prediction looks like, not just a number.
\end{showbox}

\begin{measurebox}
Each panel is one test sample chosen at an error percentile (2, 25, 50, 75, 90, 98). Color = predicted
probability over the floor; \textbf{red $\times$} = predicted peak; \textbf{white $\circ$} = true
location. The distance between $\times$ and $\circ$ is that sample's error.
\end{measurebox}

\begin{madebox}
The six samples and their heatmaps were saved by \texttt{dump\_results.py} into
\texttt{heatmaps\_unet\_fig10b.npz}; \texttt{plot\_heatmaps.py} rendered them with the peak/truth markers
and gamma-brightening.
\end{madebox}

\begin{readbox}
Low-error panels (left) show a \textbf{single tight blob} with $\times$ right on $\circ$. The high-error
panel shows a \textbf{multimodal} heatmap (several bright spots) where the model picked the wrong one ---
this is what causes the rare large errors (the tail of the CDF).
\end{readbox}

% ---------------------------------------------------------------
\subsection{Qualitative in-distribution --- MobileNetV2 (Appendix)}
\figpdf[1.0]{heatmaps_mobilenet_fig10b.pdf}

\begin{showbox}
The same six-panel view for the plain \textbf{MobileNetV2}, to compare its prediction quality against the
U-Net's.
\end{showbox}

\begin{readbox}
Predictions are good and tight (it matches DLoc at 0.707\,m), but the blobs are slightly less sharp than
the U-Net's --- consistent with the U-Net's small accuracy edge.
\end{readbox}

% ---------------------------------------------------------------
\subsection{Qualitative in-distribution --- TinyCNN (Appendix)}
\figpdf[1.0]{heatmaps_tinycnn_fig10b.pdf}

\begin{showbox}
The six-panel view for \textbf{TinyCNN} --- the model that failed (3.9\,m median).
\end{showbox}

\begin{readbox}
The heatmaps are \textbf{broad and not peaked} --- the model spreads its probability instead of
committing to a spot, so $\times$ often sits far from $\circ$. This is the visual signature of a model
\emph{below the $\sim$2M capacity floor}: too little capacity to form a sharp, confident prediction.
\end{readbox}

% =====================================================================
\section{Cross-Session Figures (does it transfer to a new environment?)}
% =====================================================================

These come from the leave-one-session-out protocol: train on July + two August sessions, test on the
\textbf{held-out} August session. They answer: \emph{``does it still work when the room changes, and how
does it compare to DLoc fairly?''}

% ---------------------------------------------------------------
\subsection{Leave-one-session-out median bars (Thesis Fig.~5.3)}
\figpdf[0.78]{crosssession_median.pdf}

\begin{showbox}
A bar chart comparing our \textbf{2.41M-parameter U-Net} against the published \textbf{16.5M-parameter
DLoc}, on each of the three held-out August environments. This is the headline ``fair comparison.''
\end{showbox}

\begin{measurebox}
Each bar's height = the \textbf{median} cross-session error (meters) on that held-out fold. Our bars are
the \textbf{three-seed mean}, with an error bar showing the standard deviation across seeds (42/123/777).
\end{measurebox}

\begin{madebox}
\texttt{plot\_crosssession.py} read the per-fold, per-seed results from
\texttt{code/results\_crosssession/*.json} (nine training runs: 3 folds $\times$ 3 seeds, 120 epochs
each), averaged the seeds, and plotted them beside the DLoc paper's published numbers.
\end{madebox}

\begin{readbox}
Our U-Net \textbf{matches DLoc on the different-furniture fold} (0.822 vs 0.82\,m) and is within
$\sim$0.1\,m on the others --- \emph{at 6.8$\times$ fewer parameters}. The \textbf{reflector fold is
clearly the tallest} (hardest), as expected from the physics. The small error bars show the result is
stable across seeds (not luck).
\end{readbox}

% ---------------------------------------------------------------
\subsection{Cross-session error CDFs (Thesis Fig.~5.4)}
\figpdf[0.78]{cdf_crosssession.pdf}

\begin{showbox}
The full error distribution (not just the median) of the U-Net on each of the three held-out folds, on
one plot.
\end{showbox}

\begin{measurebox}
x-axis: error in meters; y-axis: fraction of that fold's test samples with error $\le x$. One curve per
held-out August session (11{,}201 / 8{,}754 / 5{,}486 samples respectively).
\end{measurebox}

\begin{madebox}
From \texttt{errors\_unet\_cross\_env2/3/4.npy} (the per-sample errors on each held-out fold), plotted by
\texttt{plot\_cdfs.py}.
\end{madebox}

\begin{readbox}
The two \textbf{furniture folds track closely} together, while the \textbf{reflector fold
(Aug16\_4\_ref) is shifted to larger errors} --- its curve climbs more slowly, with a heavier tail. This
is the distribution-level view of why the reflector is the hardest environment.
\end{readbox}

% ---------------------------------------------------------------
\subsection{Qualitative cross-session --- reflector fold (Thesis Fig.~5.5)}
\figpdf[1.0]{heatmaps_unet_cross_env4.pdf}

\begin{showbox}
Six predicted heatmaps on the hardest fold (the \textbf{aluminum reflector} session), to \emph{see} why
the errors grow there.
\end{showbox}

\begin{measurebox}
Same format as the in-distribution heatmaps (red $\times$ = predicted peak, white $\circ$ = truth, six
error percentiles), but on the unseen reflector environment.
\end{measurebox}

\begin{readbox}
Compared with the clean in-distribution panels, these show \textbf{extra scattered probability mass} ---
the ``phantom'' specular reflections from the metal panel. That extra bright junk is what pulls the
predicted peak away from the truth and drives the larger errors on this fold.
\end{readbox}

% ---------------------------------------------------------------
\subsection{Qualitative cross-session --- furniture folds (Appendix)}
\figpdf[1.0]{heatmaps_unet_cross_env2.pdf}
\figpdf[1.0]{heatmaps_unet_cross_env3.pdf}

\begin{showbox}
The same six-panel view on the two \emph{furniture} folds: Aug16\_1 (furniture rearrangement) and
Aug16\_3 (different furniture).
\end{showbox}

\begin{readbox}
Predictions stay mostly tight and on-target --- noticeably cleaner than the reflector fold above. This
visually confirms the difficulty ordering: \textbf{furniture changes are mild domain shifts; the metal
reflector is the severe one.}
\end{readbox}

% =====================================================================
\section{Efficiency Figures (is it actually lighter?)}
% =====================================================================

These come from the efficiency benchmark (dummy inputs, no data, no training). They answer: \emph{``what
did we pay --- or save --- in size and speed?''}

% ---------------------------------------------------------------
\subsection{Accuracy versus model size (Thesis Fig.~5.6)}
\figpdf[0.7]{compression_pareto.pdf}

\begin{showbox}
A scatter of every model placed by \textbf{how accurate} it is and \textbf{how big} it is --- so you can
see the best trade-off at a glance.
\end{showbox}

\begin{measurebox}
One axis is median error (lower = better), the other is parameter count / model size. \textbf{The ideal
corner is the lower-left}: low error \emph{and} small size.
\end{measurebox}

\begin{madebox}
\texttt{plot\_compression.py}, using the verified accuracy and parameter numbers (DLoc 16.5M, U-Net
2.41M, MobileNetV2 2.3M, TinyCNN 47K, Mamba 0.65M).
\end{madebox}

\begin{readbox}
The \textbf{MobileNetV2-UNet sits in the favorable lower-left} --- best median error at a small fraction
of DLoc's size. Mamba sits at the top (worst error); TinyCNN is small but far to the high-error side.
This single plot is the visual summary of the whole compression story.
\end{readbox}

% ---------------------------------------------------------------
\subsection{Parameter counts (Thesis Fig.~5.7)}
\figpdf[0.7]{compression_params.pdf}

\begin{showbox}
A simple bar chart of how many parameters each model has, on a \textbf{log scale} (so the huge range from
47K to 16.5M fits on one plot).
\end{showbox}

\begin{measurebox}
Bar height = parameter count (log scale). The log axis means each equal step is a $10\times$ change, not
a fixed amount.
\end{measurebox}

\begin{readbox}
Shows the scale of compression at a glance: the U-Net and MobileNetV2 bars are far below DLoc
($\sim$7$\times$ smaller), and TinyCNN is tiny ($351\times$ smaller) --- but recall from the accuracy
figures that TinyCNN paid for that with poor accuracy.
\end{readbox}

% =====================================================================
\section{One-screen summary}
% =====================================================================

\begin{table}[h]
\centering
\footnotesize
\begin{tabular}{@{}lll@{}}
\toprule
\textbf{Figure} & \textbf{Measures} & \textbf{Main message} \\
\midrule
cdf\_fig10b              & in-dist error distribution      & U-Net best across the whole range \\
heatmaps\_unet\_fig10b   & in-dist predicted heatmaps      & tight blobs; tail = multimodal cases \\
heatmaps\_mobilenet      & in-dist predicted heatmaps      & good, slightly less sharp than U-Net \\
heatmaps\_tinycnn        & in-dist predicted heatmaps      & broad/unpeaked = below capacity floor \\
crosssession\_median     & LOSO median vs DLoc             & competitive at 6.8$\times$ smaller \\
cdf\_crosssession        & cross-session error distrib.\   & reflector fold hardest (heavy tail) \\
heatmaps\_unet\_cross\_env4 & reflector-fold heatmaps      & phantom reflections cause errors \\
heatmaps\_unet\_cross\_env2/3 & furniture-fold heatmaps   & mild shift, predictions stay tight \\
compression\_pareto      & accuracy vs size               & U-Net at the ideal lower-left \\
compression\_params      & parameter counts (log)         & $\sim$7$\times$ smaller than DLoc \\
\bottomrule
\end{tabular}
\end{table}

\begin{readbox}[The figures tell one story]
The \textbf{CDFs and heatmaps} prove the small U-Net is as accurate as DLoc (and show \emph{why} the hard
cases are hard --- multipath and reflections); the \textbf{cross-session bars and CDFs} prove it
\emph{transfers} to new rooms; and the \textbf{efficiency plots} prove it does all this at a fraction of
the size and compute.
\end{readbox}

\end{document}
